\chapter{Ptosis Repair}

\section*{What Is This Surgery?}
Ptosis repair is a specialized surgical procedure designed to correct blepharoptosis, commonly known as a drooping upper eyelid. This condition can affect one or both eyes, occurring congenitally due to developmental muscle dystrophy or developing later in life from various acquired causes, most commonly aponeurotic dehiscence. The primary goal of the surgery is to elevate the upper eyelid to a more normal, functional, and aesthetically pleasing position, thereby improving the patient's superior visual field and often addressing significant cosmetic concerns. By restoring proper eyelid height and contour, ptosis repair aims to alleviate symptoms such as visual obstruction, compensatory brow ache, and fatigue, ultimately enhancing both visual function and overall quality of life. The choice of surgical technique is highly individualized, depending on the underlying etiology of the ptosis and the functional status of the levator palpebrae superioris muscle.

\section*{Alternative Names}
\begin{itemize}
  \item Blepharoptosis Repair
  \item Eyelid Lift Surgery (though distinct from blepharoplasty)
  \item Levator Resection or Advancement
  \item Frontalis Sling Procedure
  \item M\"uller's Muscle Conjunctival Resection (MMCR)
  \item Aponeurotic Ptosis Repair
  \item Eyelid Droop Correction
\end{itemize}

\section*{Relevant Surgical Anatomy}
A thorough understanding of the intricate anatomy of the upper eyelid and surrounding structures is paramount for successful ptosis repair. The upper eyelid is a complex structure composed of several layers, from superficial to deep: skin, orbicularis oculi muscle, orbital septum, preaponeurotic fat pads, levator palpebrae superioris muscle and its aponeurosis, M\"uller's muscle, and the conjunctiva lining the posterior surface of the tarsal plate. The tarsal plate, a dense fibrous structure, provides the main structural support for the eyelid.

The primary elevator of the upper eyelid is the levator palpebrae superioris muscle, originating from the lesser wing of the sphenoid bone and innervated by the superior division of the oculomotor nerve (cranial nerve III). It transitions into a broad aponeurosis that inserts onto the anterior surface of the tarsal plate and into the skin to form the eyelid crease. Deep to the levator aponeurosis lies M\"uller's muscle, a sympathetically innervated smooth muscle that contributes 1-2 mm of eyelid elevation. The orbicularis oculi muscle, innervated by the facial nerve (cranial nerve VII), is responsible for eyelid closure. Arterial supply to the eyelids is primarily from the medial and lateral palpebral arteries, branches of the ophthalmic artery, forming marginal and peripheral arcades. Venous drainage parallels the arterial supply, emptying into the ophthalmic and facial veins. Sensory innervation is provided by branches of the trigeminal nerve (cranial nerve V). Lymphatic drainage from the lateral two-thirds of the upper eyelid drains to the preauricular lymph nodes, while the medial one-third drains to the submandibular nodes. Key surgical landmarks include the superior limbus, the eyelid crease, the tarsal plate, and the orbital septum, which separates the preaponeurotic fat from the orbital contents.

\section{Surgical Principle \& Rationale}
The fundamental principle of ptosis repair is to restore the upper eyelid margin to a physiologically appropriate position, typically 1-2 mm below the superior limbus, ensuring clear vision and aesthetic symmetry. The specific surgical technique employed depends critically on the underlying cause of the ptosis, the severity of the droop, and the functional status of the levator palpebrae superioris muscle, which is assessed preoperatively by measuring levator function (excursion). The primary mechanisms involve either strengthening or shortening the muscles responsible for eyelid elevation, or utilizing an alternative muscle group to lift the eyelid.

For aponeurotic ptosis, the most common type in adults, the procedure focuses on repairing the stretched, dehisced, or disinserted levator aponeurosis. This often involves advancing, resecting, or plicating the aponeurosis and reattaching it to the anterior surface of the tarsal plate, effectively tightening the eyelid's primary elevator. In cases of mild ptosis (typically 1-2 mm) with good M\"uller's muscle function and a positive phenylephrine test, a transconjunctival M\"uller's muscle conjunctival resection (MMCR) may be performed, relying on the sympathetic innervation of this muscle. For severe ptosis with poor levator function (typically less than 4-5 mm of excursion), particularly in children with congenital ptosis, a frontalis suspension (sling) procedure is often necessary. This technique connects the tarsal plate to the frontalis muscle, allowing the eyebrow to lift the eyelid. The overarching technical goals are to achieve stable, symmetric eyelid height, a natural contour, and full visual clearance, while minimizing the risk of lagophthalmos (incomplete eyelid closure), exposure keratopathy, and dry eye.

\section{History \& Surgical Evolution}
The surgical correction of ptosis has evolved significantly since its early attempts in the late 19th century. Initial procedures were often rudimentary, focusing on simple skin excisions or muscle plications. Dransart is credited with early descriptions of levator plication in the 1880s. In the early 20th century, more anatomically precise techniques began to emerge. Blaskovics, in 1923, described an internal levator resection technique that laid the groundwork for modern approaches by directly addressing the levator muscle. Concurrently, surgeons like Everbusch (1890s) and Friedenwald (1930s) explored methods of frontalis suspension, utilizing the brow muscle to elevate the eyelid in cases of severe levator weakness, often using sutures or strips of skin.

A significant advancement came in 1961 when Fasanella and Servat introduced the M\"uller's muscle conjunctival resection, a transconjunctival technique still widely used for mild ptosis with good M\"uller's muscle function. Over the latter half of the 20th century, surgeons such as Putterman, Beard, and Jones further refined external levator advancement and resection techniques, emphasizing precise anatomical repair, intraoperative adjustment of eyelid height (especially under local anesthesia), and the importance of recreating a natural eyelid crease. The development of various sling materials for frontalis suspension, from autogenous fascia lata (first described by Wright in 1922) to synthetic materials like silicone rods, also marked progress, offering durable solutions for complex cases of poor levator function. Modern techniques continue to focus on individualized approaches, minimizing complications, and optimizing both functional and aesthetic outcomes.

\section{Clinical Indications}
Ptosis repair is indicated for a variety of clinical scenarios, primarily when the drooping eyelid causes functional impairment or significant cosmetic distress.

\begin{itemize}
  \item Significant obstruction of the superior visual field, leading to difficulty with reading, driving, computer use, or other daily activities, often quantifiable by visual field testing.
  \item Chronic brow ache, headache, or fatigue resulting from compensatory over-activation of the frontalis muscle to lift the eyelid, a common symptom of acquired ptosis.
  \item Amblyopia (lazy eye) in children with congenital ptosis, where the drooping eyelid obstructs the visual axis during critical developmental periods, necessitating early intervention.
  \item Cosmetic concerns impacting the patient's quality of life, even in the absence of severe visual obstruction, due to perceived asymmetry or an aged appearance.
  \item Asymmetry between the eyelids that causes functional or aesthetic distress, often requiring bilateral surgery to achieve optimal symmetry.
  \item Acquired ptosis secondary to aponeurotic dehiscence (most common), neurogenic causes (e.g., partial oculomotor nerve palsy, Horner's syndrome after appropriate workup), or myogenic causes (e.g., myasthenia gravis, chronic progressive external ophthalmoplegia, after medical management fails).
  \item Pseudoptosis due to dermatochalasis (excess eyelid skin) or brow ptosis, where the eyelid position itself is normal but appears low, often requiring combined procedures.
  \item Traumatic ptosis where the levator complex has been damaged, requiring reconstruction.
\end{itemize}

\section{Clinical Positioning}
Ptosis repair is predominantly a primary, definitive surgical treatment for symptomatic blepharoptosis. For most patients presenting with acquired aponeurotic ptosis, which is the most common form in adults, or congenital ptosis, it is the initial and often curative intervention aimed at restoring both visual function and aesthetic balance. The choice of specific technique (e.g., external levator advancement, M\"uller's muscle conjunctival resection, or frontalis suspension) is determined by the severity of ptosis and the underlying levator function, making it a primary procedure with varying approaches tailored to the individual clinical presentation.

However, ptosis repair can also serve as a secondary or revision procedure. This is common in cases where an initial repair was unsuccessful, resulting in under-correction (persistent ptosis), over-correction (eyelid retraction), recurrence of the ptosis, or unsatisfactory contour. In such scenarios, the goal of the secondary surgery is to refine the eyelid position and contour, often requiring more complex dissection due to scar tissue. While not typically an adjunct to other non-eyelid procedures, it may be performed concurrently with other oculoplastic surgeries such as blepharoplasty (for dermatochalasis) or brow lift (for brow ptosis) to achieve a comprehensive periorbital rejuvenation.

\section{Surgical Technique \& Procedure}
Ptosis repair is typically performed under local anesthesia with intravenous sedation, allowing for intraoperative patient cooperation to assess eyelid height and symmetry. General anesthesia may be used for children, anxious adults, or complex revision cases. The patient is positioned supine with the head slightly elevated.

The procedure begins with precise preoperative markings of the desired eyelid crease, the amount of skin to be excised (if combined with blepharoplasty), and the target eyelid height relative to the superior limbus. Local anesthetic (e.g., lidocaine with epinephrine) is then infiltrated into the upper eyelid.

The most common approach for aponeurotic ptosis is the external levator advancement/resection:
\begin{enumerate}
  \item  **Incision:** An incision is made along the natural upper eyelid crease, extending laterally.
  \item  **Dissection:** The skin and underlying orbicularis oculi muscle are carefully dissected to expose the orbital septum.
  3.  **Identification of Structures:** The orbital septum is opened, revealing the preaponeurotic fat pads. These are gently retracted superiorly to expose the underlying levator aponeurosis.
  4.  **Aponeurosis Repair:** The surgeon identifies the stretched, dehisced, or disinserted levator aponeurosis. The aponeurosis is then advanced, resected, or plicated to shorten its effective length.
  5.  **Tarsal Fixation \& Adjustment:** The repositioned aponeurosis is reattached to the anterior surface of the tarsal plate with fine, non-absorbable or slowly absorbable sutures. With the patient awake (if under local anesthesia), eyelid height and contour are meticulously adjusted until optimal symmetry and position are achieved.
  6.  **Crease Formation:** The skin and orbicularis muscle are then sutured to the newly positioned aponeurosis to recreate a natural eyelid crease.
  7.  **Closure:** The skin incision is closed with fine sutures, which may be absorbable or non-absorbable.
\end{enumerate}

For M\"uller's muscle conjunctival resection (MMCR), a transconjunctival approach is used, where the M\"uller's muscle and a portion of the conjunctiva are resected from the undersurface of the eyelid. In cases of severe levator dysfunction, a frontalis suspension procedure is performed, where a sling material (e.g., autogenous fascia lata, banked fascia lata, or a silicone rod) is passed from the tarsal plate, through the orbicularis and orbital septum, to be anchored to the frontalis muscle in the brow, allowing the brow to lift the eyelid. Drains are typically not required, and implants are only used in frontalis sling procedures.

\section*{Preoperative Workup \& Preparation}
Thorough preoperative preparation is crucial for a successful ptosis repair, focusing on patient safety and optimizing surgical outcomes. This begins with a comprehensive ophthalmologic examination, including detailed history, visual acuity, visual field testing (especially superior field), ocular motility assessment, tear film evaluation (e.g., Schirmer's test, tear break-up time), measurement of levator function (excursion), margin reflex distance (MRD1 and MRD2), and assessment for Bell's phenomenon and lagophthalmos. A detailed medical history and physical examination are conducted to identify any comorbidities, allergies, and current medications, with particular attention to anticoagulants or antiplatelet agents.

\begin{itemize}
  \item **Pre-anesthesia evaluation:** Routine laboratory tests (complete blood count, electrolytes, coagulation profile), electrocardiogram, and chest X-ray may be ordered based on the patient's age, medical history, and planned anesthesia type.
  \item **NPO status:** Patients are typically instructed to be NPO (nothing by mouth) for 6-8 hours prior to surgery to prevent aspiration risks associated with anesthesia.
  \item **Medication adjustments:** Blood-thinning medications such as aspirin, NSAIDs, and warfarin are usually discontinued 7-10 days pre-operatively, under the guidance of the prescribing physician, to minimize bleeding risk. Other medications may be adjusted as per anesthesiologist's recommendations.
  \item **Prophylactic antibiotics:** Intravenous antibiotics (e.g., cefazolin) are typically administered just prior to the incision to reduce the risk of surgical site infection. Topical antibiotic drops or ointment may also be prescribed for a few days pre-op.
  \item **Informed consent:** A detailed discussion of the procedure, its potential benefits, risks, alternatives, and the possibility of revision surgery is conducted, and the patient provides informed consent. Preoperative photographs are routinely taken for documentation and comparison.
\end{itemize}

\section*{Contraindications \& Precautions}
Careful patient selection is essential to minimize risks and optimize outcomes for ptosis repair.

Absolute Contraindications:
\begin{itemize}
  \item Uncontrolled severe dry eye syndrome, as surgery can exacerbate symptoms and lead to sight-threatening exposure keratopathy.
  \item Active ocular infection (e.g., conjunctivitis, blepharitis, hordeolum) in the operative eye, which must be treated prior to surgery.
  \item Unrealistic patient expectations regarding the surgical outcome, which can lead to dissatisfaction despite a technically successful procedure.
  \item Severe underlying medical conditions that preclude safe administration of anesthesia or pose an unacceptable surgical risk (e.g., unstable angina, recent myocardial infarction).
  \item Untreated thyroid eye disease in its active inflammatory phase, as eyelid position can fluctuate.
\end{itemize}

Relative Contraindications and Precautions:
\begin{itemize}
  \item Insufficient levator function for specific techniques (e.g., less than 4 mm excursion for levator resection), requiring consideration of alternative procedures like frontalis suspension.
  \item Pre-existing significant lagophthalmos or poor Bell's phenomenon, which increases the risk of postoperative exposure keratopathy and necessitates a conservative surgical approach.
  \item Bleeding diathesis or current use of anticoagulant medications, requiring careful management, temporary cessation, or bridging therapy under medical supervision.
  \item Unstable or active dermatological conditions affecting the eyelids (e.g., severe eczema, psoriasis) that could impair wound healing.
  \item Certain progressive neurological conditions (e.g., myasthenia gravis not adequately controlled, chronic progressive external ophthalmoplegia) that may affect future eyelid function and require careful consideration of long-term stability.
  \item Monocular patients, where the risk of complications to the sole seeing eye must be weighed heavily against potential benefits.
\end{itemize}

\section*{Complications \& Risks}
As with any surgical procedure, ptosis repair carries inherent risks, which can be broadly categorized as general surgical risks and procedure-specific complications.

General Surgical Risks:
\begin{itemize}
  \item **Bleeding:** May manifest as ecchymosis (bruising), hematoma formation, or, rarely, retrobulbar hemorrhage, which is an ophthalmic emergency that can cause optic nerve compression and threaten vision.
  \item **Infection:** Cellulitis of the eyelid, wound infection, or abscess formation, requiring antibiotic treatment and potentially surgical drainage.
  \item **Anesthesia complications:** Nausea, vomiting, allergic reactions to anesthetic agents, or more serious cardiac or respiratory events, particularly with general anesthesia.
  \item **Pain:** Postoperative discomfort, usually mild and manageable with oral analgesics.
\end{itemize}

Procedure-Specific Risks:
\begin{itemize}
  \item **Asymmetry of eyelid height or contour:** The most common complication, which may necessitate revision surgery to achieve optimal balance.
  \item **Under-correction:** Persistent ptosis, where the eyelid remains too low, requiring further surgical intervention.
  \item **Over-correction:** Eyelid retraction or lagophthalmos (incomplete eyelid closure), leading to exposure keratopathy, dry eye, and potential corneal damage if severe.
  \item **Dry eye syndrome:** Exacerbation of pre-existing dry eye or new onset symptoms due to altered tear film dynamics and increased ocular surface exposure.
  \item **Suture granuloma or cyst formation:** Along the incision line or within the eyelid, sometimes requiring removal.
  \item **Scarring:** Hypertrophic scarring or keloid formation, though rare in the eyelids, can occur.
  \item **Eyelid contour abnormalities:** Such as a peaked, flattened, or notched arch, affecting aesthetic outcome.
  \item **Damage to the globe:** Rare but serious complications include corneal abrasion, globe perforation, or vision loss due to direct trauma during surgery.
  \item **Diplopia (double vision):** May occur due to injury to the superior rectus muscle or its innervation, particularly with aggressive dissection during levator advancement.
  \item **Loss of eyelid crease:** Or an unnatural-looking crease.
  \item **Recurrence of ptosis:** Over time, especially with progressive underlying conditions or age-related changes.
\end{itemize}

\section*{Postoperative Recovery Timeline}
Immediately following ptosis repair, the patient is transferred to the Post-Anesthesia Care Unit (PACU) for monitoring of vital signs and initial recovery from anesthesia. Ice packs are typically applied to the eyelids intermittently to minimize swelling and bruising, and head elevation is encouraged to reduce edema. Mild discomfort is common and usually managed with oral analgesics. Vision is checked, and ophthalmic antibiotic ointment is applied to the incision sites. Most patients are discharged home the same day with detailed post-operative instructions.

During the first 24-48 hours, continued intermittent application of ice packs (15-20 minutes on, 20-30 minutes off) is highly recommended. Swelling, bruising (ecchymosis), and mild discomfort are expected and gradually subside. Patients are advised to avoid strenuous activities, heavy lifting, bending over, and rubbing their eyes. Ophthalmic antibiotic ointment should be applied to the incision as directed, typically twice daily. Over the first 1-2 weeks, swelling and bruising continue to resolve, though some residual puffiness may persist. Non-absorbable skin sutures, if used, are typically removed around 5-7 days post-operatively. Patients should refrain from wearing eye makeup or using contact lenses until cleared by the surgeon, usually after 2-3 weeks. Most patients can return to light activities and work within 1-2 weeks. Full resolution of swelling and the final aesthetic outcome may take several weeks to months (3-6 months) as tissues heal and remodel.

\section*{Surgical Findings \& Pathology}
During ptosis repair, the surgeon meticulously documents several key findings that guide the procedure and inform the postoperative assessment. These include the specific type of ptosis (e.g., aponeurotic, congenital, mechanical), the measured levator function, the degree of aponeurotic dehiscence, stretching, or disinsertion, the presence and amount of preaponeurotic fat, and the final intraoperative eyelid height and contour achieved. For congenital ptosis, the levator muscle itself may appear thin, fibrotic, or infiltrated with fat. The amount of levator aponeurosis or M\"uller's muscle resected or advanced is also recorded.

Pathological findings are generally not relevant unless an underlying mass or lesion was biopsied or excised during the procedure. In the case of aponeurotic ptosis, the resected tissue, if sent for pathology, typically confirms fibrous tissue consistent with levator aponeurosis. For congenital ptosis, a levator muscle biopsy might reveal a dysgenetic muscle with varying degrees of fatty infiltration or fibrosis, confirming the diagnosis of a developmental dystrophy. In rare instances where a tumor or inflammatory lesion is suspected, the excised tissue would be sent for histopathological examination to establish a definitive diagnosis.

\section{Expected Postoperative Course}
A normal and successful postoperative course following ptosis repair is characterized by a gradual and progressive improvement in eyelid appearance and function. Patients can expect initial swelling and bruising to resolve over 1-2 weeks, with mild discomfort that is well-controlled by oral analgesics. The eyelid margin should be positioned appropriately, typically 1-2 mm below the superior limbus, with good symmetry between both eyes if bilateral surgery was performed. The patient should experience a significant improvement in their superior visual field, leading to enhanced functional vision and a reduction in symptoms such as chronic brow ache or fatigue. The eyelid contour should appear natural and aesthetically pleasing, without significant peaking or flattening. Ideally, there should be no significant lagophthalmos (incomplete eyelid closure), new or exacerbated dry eye symptoms, or visual disturbances. The incision lines should heal well, becoming inconspicuous over time, blending into the natural eyelid crease within several months. Patients typically report high satisfaction with both the functional and cosmetic improvements.

\section{Postoperative Care \& Follow-up}
Postoperative care involves a structured follow-up schedule to monitor healing, assess outcomes, and address any potential complications.

\begin{itemize}
  \item **First post-operative visit:** Typically scheduled 5-7 days after surgery for suture removal (if non-absorbable) and initial assessment of wound healing, eyelid position, and ocular surface health.
  \item **Subsequent visits:** Usually at 1 month, 3 months, and 6-12 months to monitor long-term results, assess for symmetry, stability of eyelid height, and address any evolving concerns or the need for revision.
  \item **Medications:** Patients are typically prescribed topical antibiotic ointment for 1-2 weeks and may use artificial tears or lubricating gels to manage dry eye symptoms. Oral analgesics are used as needed for pain.
  \item **Activity resumption:** Patients are advised to gradually return to normal activities, avoiding strenuous exercise, heavy lifting, and swimming for 2-4 weeks. Contact sports should be avoided for a longer period.
  \item **Warning signs:** Patients are instructed to contact their surgeon immediately for sudden vision changes, severe or worsening pain, excessive swelling, redness, purulent discharge from the wound, fever, or persistent inability to close the eye.
\end{itemize}
Revision surgery may be considered for significant under-correction, over-correction, or persistent asymmetry, but is typically deferred until adequate healing time has passed, which is usually 6-12 months post-initial surgery, to allow for complete resolution of swelling and scar maturation.

\section*{Epidemiology}
Ptosis affects individuals across all age groups, exhibiting a bimodal distribution. Congenital ptosis, present at birth, often results from a developmental dystrophy of the levator muscle and is frequently unilateral. Its prevalence is estimated to be around 1 in 840 live births, and if severe, it can lead to amblyopia (lazy eye) if not corrected early. Acquired ptosis is far more common, with aponeurotic ptosis being the most frequent type, accounting for over 70\% of adult cases. This age-related condition typically affects older adults due to the dehiscence, disinsertion, or stretching of the levator aponeurosis, often exacerbated by chronic eye rubbing or contact lens wear. Neurogenic ptosis (e.g., Horner's syndrome, third nerve palsy) and myogenic ptosis (e.g., myasthenia gravis, chronic progressive external ophthalmoplegia) are less common but require specific diagnostic workups. While precise population-based incidence figures for ptosis requiring surgical correction vary, it represents a significant portion of oculoplastic surgical volume globally. Mortality directly related to ptosis repair is exceedingly rare, primarily associated with general anesthesia risks in medically fragile patients. Morbidity is usually related to aesthetic outcomes, minor functional issues like dry eye, or the need for revision surgery.

\section*{Outcomes \& Prognosis}
The outcomes of ptosis repair are generally favorable, with high success rates in achieving satisfactory eyelid height, contour, and visual improvement. For common aponeurotic ptosis repairs, success rates, defined as achieving an eyelid margin within 1 mm of the desired position with good contour, range from 85\% to 95\%. Functional outcomes include a significantly improved superior visual field, reduced compensatory brow fatigue, and an overall enhancement in the patient's quality of life. Ptosis repair does not directly impact survival outcomes. Recurrence of ptosis can occur, particularly in cases of severe levator dysfunction, progressive underlying conditions, or inadequate initial correction, necessitating revision surgery in approximately 5-15\% of patients.

Prognostic factors for a good outcome include accurate preoperative assessment, adequate residual levator function (for levator-based repairs), absence of severe dry eye, and patient compliance with post-operative care instructions. For congenital ptosis, early intervention is crucial to prevent amblyopia. While long-term stability is generally good, age-related changes can lead to a gradual recurrence of ptosis over many years, potentially requiring further intervention. Patient satisfaction is generally high, especially when expectations are appropriately managed preoperatively.

\section*{References}
\begin{enumerate}
  \item MedlinePlus. Ptosis surgery. U.S. National Library of Medicine; 2024. Available from: \url{https://medlineplus.gov/ency/article/002986.htm}
  \item Kanski JJ, Bowling B. Kanski's Clinical Ophthalmology: A Systematic Approach. 9th ed. Elsevier; 2020.
  \item American Academy of Ophthalmology. Preferred Practice Pattern: Adult Ptosis. San Francisco, CA: American Academy of Ophthalmology; 2019.
  \item Yanoff M, Duker JS. Ophthalmology. 5th ed. Elsevier; 2019.
  \item Putterman AM. Putterman's Cosmetic Oculoplastic Surgery. 4th ed. Saunders; 2008.
\end{enumerate}

\clearpage